% 本文件由 LaTeX 排版 Agent 根据有效 translation.json 自动生成。
% author=Apocrypha; work_id=0805; title=若瑟木匠史

\chapter{\WorkTitle{若瑟木匠史}}

% structure: p0001 source=h1 target=omitted reason=page-title-already-rendered-by-parent

因一体三位的天主之名。\par

我们的父亲、圣善老人、木匠若瑟逝世史。\par

愿他的祝福与祈祷护守我们众人，诸位弟兄！阿们。\par

他的一生共一百一十一年，离世是在阿彼布月（Abib）二十六日，即阿布月（Ab）。愿他的祈祷护守我们！阿们。事实上，是我们的主耶稣基督亲自在\Place{橄榄山}上将这段历史、若瑟的一切劳苦以及他日子的终结告诉了祂的圣门徒。圣宗徒们保存了这次谈话，并记录下来，存放在耶路撒冷的图书馆。愿他们的祈祷护守我们！阿们。\par

一日，救主、我们的师尊、天主和救主耶稣基督与祂的门徒同坐，他们全聚集在\Place{橄榄山}上，祂对他们说：我的弟兄和朋友们，圣父从众人中拣选的子们，你们知道我曾多次告诉你们，我必须被钉十字架，必须为亚当及其后裔的救恩而死，并且我要从死者中复活。现在我要将先前向你们宣布的圣福音的道理托付给你们，使你们能向全世界宣讲。我要从高处赐给你们能力，并以圣神充满你们。\BibleRef{路 24:49} 并要你们向万民宣告悔改与罪过的赦免。\BibleRef{路 24:37} 因为一杯水，\BibleRef{玛 10:42} 若有人在来世获得它，这杯水就比这整个世界的所有财富更大、更好。在我父家里，一足之地也比地上的一切财富更大、更卓越。是的，虔诚者的喜乐居所中的一个时辰，比在罪人中的一千年更蒙福、更珍贵：因为他们的哭泣和哀叹不会终止，他们的眼泪不会停止，他们也永远找不到安慰和安息。现在，我尊贵的肢体们，去向万民宣告，告诉他们，对他们说：救主勤勉查考应得的产业，并是正义的管理者。天使们将打倒他们的敌人，并在争战之日为他们争战。祂要审察人所说的每一句愚妄和闲话，并要为此交账。\BibleRef{玛 12:36} 因为没有人能逃脱死亡，同样，每个人的作为在审判之日也必显露出来，或善或恶。\BibleRef{格后 5:10} 也告诉他们我今天对你们所说的这句话：愿强壮者不要夸耀自己的强壮，富足者不要夸耀自己的财富；但谁愿夸耀，就当在主内夸耀。\par

有一个人名叫\Person{若瑟}，出生于犹大地的\Place{白冷}城，即达味王的城。此人富有智慧和学识，曾在上主圣殿中担任司祭。此外，他擅长手艺，是一名木匠；如同众人一样，他娶了妻子。他生了儿女，四个儿子和两个女儿。他们的名字是：\Person{犹达斯}、\Person{犹斯托}、\Person{雅各伯}、\Person{西满}。两个女儿的名字是\Person{阿西亚}和\Person{里狄雅}。后来，义人若瑟的妻子——一位在所有工作中都关注天主光荣的妇人——离世了。而若瑟，这位义人，我肉身之父和我母亲玛利亚的配偶，便带着他的儿子们去从事他的手艺，操练木匠之业。\par

当义人若瑟成为鳏夫时，我的母亲\Person{玛利亚}——蒙福、圣洁、纯洁——已经十二岁了。她的父母在她三岁时将她献于圣殿，她在上主的圣殿中度过了九年。当司祭们看到这位圣洁而敬畏天主的贞女渐渐长大，便彼此商议说：让我们寻找一位义人且虔诚的人，将玛利亚托付给他，直到她结婚的时候；免得她继续留在圣殿中，发生妇女常有事，以致我们犯罪，天主向我们发怒。\par

于是他们立即派人去，召集了犹大支派的十二位老人，记下了以色列十二支派的名字。抽签后，签落在这位虔诚的老人、义人若瑟身上。司祭们便对我的蒙福母亲说：与若瑟同去，与他同住，直到你结婚的时候。义人若瑟于是接了我的母亲，领她到自己的家。玛利亚在他父亲家里发现了小雅各伯，因丧母而心碎悲伤，便抚养他。因此玛利亚被称为雅各伯的母亲。\BibleRef{路 24:10} 此后，若瑟留她在家里，自己去了他做木匠手艺的作坊。圣贞女在他家住满两年后，她的年龄恰好十四岁，包括他接纳她的时间。\par

我出于自择，得到我圣父的同意与圣神的参议，拣选了她。我藉着一个超越受造理智理解的奥秘，从她取得了肉身。她受孕三个月后，义人若瑟从他的工作地返回；当他发现我的童贞母亲怀孕时，非常困惑，想暗暗地休了她。\BibleRef{玛 1:19} 但由于恐惧、悲伤和心中的痛苦，他那天既不能吃也不能喝。\par

但在正午，天使之长、圣\Person{加俾额尔}带着我圣父的命令，在梦中显现给他说：若瑟，达味之子，不要怕娶玛利亚为你的妻子；因为她因圣神怀孕，她要生一个儿子，将称祂为耶稣。祂要用铁杖治理万民。天使说完这话，便离开了他。若瑟从睡眠中醒来，照上主的天使对他所说的做了；玛利亚便与他同住。\BibleRef{玛 1:20}\par

7. 此后不久，奥古斯都·凯撒王颁发谕令，要天下万民各归本城登记户籍。于是义人长者若瑟起身，带着童贞玛利亚来到白冷，因为她分娩的日期已近。若瑟便将自己的名字登录在册；因为若瑟是达味的后裔，玛利亚是他的配偶，属于犹大支派。果然，我的母亲玛利亚在白冷，靠近圣祖雅各伯之妻、若瑟与本雅明之母辣黑耳的坟墓旁的一个山洞里生下了我。\par

8. 但撒殚却去把这事告诉了黑落德大帝，即阿赫劳斯的父亲。正是这个黑落德下令斩了我的朋友和亲戚若翰。因此他竭力搜寻我，以为我的王国属于这世界。\BibleRef{若 18:36}但那虔诚的长者若瑟在梦中得到警告。于是他起身，带着我的母亲玛利亚，而我躺在她的怀中。撒罗默也同他们一路同行。他们便离家出发，退往埃及，在那里住了整整一年，直到黑落德的仇恨消逝。\par

9. 黑落德以一种最悲惨的方式死去，为他所恶毒残杀的无辜婴孩之血付出代价。那邪恶的暴君黑落德死后，他们回到以色列地，住在加里肋亚一座名叫纳匝肋的城里。若瑟重操木匠旧业，亲手劳作维持生计；因为梅瑟法律吩咐，他从不指望不劳而获。\BibleRef{创 3:19}\par

10. 随着年岁增长，这位长者达到了极高的寿数。然而，他并未遭受任何身体上的衰弱，视力也未减退，口中没有掉一颗牙齿。终其一生，他的心志也从未迷乱；反而像少年人一样，在一切事务中表现出青春的活力，四肢健全，毫无疼痛。他的寿数总共一百一十一岁，高年延至极限。\par

11. 若瑟的长子犹斯托和西默盎都已成婚，各自有了家室。两个女儿也已出嫁，住在自己的家里。于是若瑟家中只剩下犹达和次雅各伯，以及我的童贞母亲。我也与他们同住，仿佛是他的儿子之一。但我一生无过。我称玛利亚为母亲，称若瑟为父亲，凡事听从他们的吩咐；从未违抗他们，而是遵行他们的命令，如同世上其他常人所做的一样；我从未激怒他们，或说出任何反对他们的话语。相反，我以极大的爱珍爱他们，如同眼中的瞳人。\par

12. 此后，那位虔诚的长者若瑟的死亡和他离世的日子临近了，正如其他源于尘土的人一样。当他的身体濒临瓦解时，主的一位天使告诉他，他的死期已经近在咫尺。因此恐惧和极大的困惑临到他。于是他起身前往耶路撒冷，进入主的圣殿，在圣所前倾心祈祷说：\par

13. 天主！一切安慰的根源，充满慈悲的天主，全人类的君主；我灵魂、身体和灵魂的天主！我以恳求敬拜祢，我主我的天主。若我的日子已尽，离世的时候临近，求祢派遣祢的圣天使长弥额尔来陪伴我；让他与我同在，使我可怜的灵魂能毫无困扰、毫无恐惧和不耐地离开这受苦的身体。因为在死亡之日，巨大的恐惧和强烈的忧伤抓住所有身体，无论男人女人，野兽或驯畜，或地上爬行的、空中飞行的。最后，天下凡有气息的活物都惊恐万状，他们的灵魂带着极大的恐惧和深重的沮丧离开身体。因此，我主我的天主，求祢让祢的圣天使以他的帮助临于我的灵魂和身体，直到它们彼此分离。并且，求祢不要使那从我出生之日就任命守护我的天使转脸离开我；愿他作我旅途的伴侣，直到带我到祢面前；愿他的面容对我欢乐愉快，愿他平安地陪伴我。也求祢不让可怖的魔鬼在我将行的路上接近我，直到我幸福地来到祢面前。并且不要让守门者阻挡我的灵魂进入乐园。求祢不要揭露我的罪过，使我在祢可畏的审判台前受谴责。不要让狮子猛扑向我；也不要让火海的波涛淹没我的灵魂——因为每个灵魂都必须经过它——直到我看见祢神性的荣耀。天主，最公义的审判者，祢以公义和公平审判世人，按各人的行为施报，我主我的天主，求祢以慈悲临于我，照亮我的道路，使我得以到达祢面前；因为祢是涌流一切美善的泉源，并永远受光荣。阿们。\par

14. 此后，当他回到纳匝肋城自己的家中时，疾病侵袭了他，他不得不卧床不起。就在这时，他按照全人类的命运去世了。因为这病对他十分沉重，他自出生以来从未像现在这样病过。的确，基督乐意如此安排义人若瑟的命运。他独身生活了四十年；此后他的妻子在他照顾下生活了四十九年，然后去世了。在他妻子去世一年后，司祭们将我的母亲、蒙福的玛利亚托付给他，要他在她婚期之前照顾她。她在若瑟家中住了两年；在她留住若瑟家的第三年，即她十五岁时，以一种受造物无人能参透或理解的奥秘——除了我自己、我的父和圣神（与我同为一体）之外——她将我生在地上。\par

15. 因此，我父亲——那位义人老者的整个寿数是一百一十一岁，这是我天上的父所预定的。他的灵魂离开身体的那日是阿彼布月二十六日。因为如今精金开始失去光辉，银子因使用而磨损——我指的是他的理智和智慧。他也厌恶饮食，丧失了他木匠手艺的一切技巧，再也不留意它了。于是在阿彼布月二十六日的黎明时分，那位义人老者若瑟躺在床上，正将他不安的灵魂交出。因此他叹息着张开嘴，双手相击，大声呼喊，说了以下的话：\par

\Quote{16. 祸哉，我生在世上的那日！祸哉，怀我的子宫！祸哉，接纳我的内脏！祸哉，哺育我的乳房！祸哉，我坐卧其上的双脚！祸哉，抚养我直到长大的双手！因为我是在不义中受孕，我母亲渴望我是在罪恶中。祸哉，我的舌头和嘴唇，它们说出了虚妄、毁谤、谎言、愚昧、嘲笑、闲话、诡诈和伪善！祸哉，我的眼睛，它们看见了可耻的事！祸哉，我的耳朵，它们喜爱毁谤者的言语！祸哉，我的双手，它们抓取了不属于自己的东西！祸哉，我的肚腹和肠子，它们贪求不可吃的食物！祸哉，我的喉咙，它像火一样吞尽一切所遇！祸哉，我的双脚，它们多次行走在天主不喜悦的道路上！祸哉，我的身体；祸哉，我可怜的灵魂，它早已背离了天主它的造主！当我到达那必须站在最公义的审判者面前的地方，当祂要我为我在青年时期积攒的行为交账时，我该做什么？祸哉，每个死在罪恶中的人！的确，那曾临到我父亲雅各伯的同样可怕时辰 \BibleRef{玛 1:16}，当他的灵魂正从身体飞出时，如今看，也临近我了。哦！我这日何等悲惨，堪当哀悼！但唯有天主是我灵魂和身体的主宰；祂也将照祂自己的美意处置它们。}\par

17. 以上是那位义人老者若瑟所说的话。而我走进他身边，发现他的灵魂极其不安，因为他陷入了极大的困惑。我对他说：\Quote{万福！我的父亲若瑟，您这位义人，您如何？}他回答我说：\Quote{万福！我亲爱的儿子。的确，死亡的痛苦和恐惧已经环绕我；但我一听到您的声音，我的灵魂就安息了。哦，纳匝肋的耶稣！耶稣，我的救主！耶稣，我灵魂的拯救者！耶稣，我的保护者！耶稣！哦，在我口中以及在一切爱慕这圣名者口中最为甘甜的名字！哦，那观看的眼睛和聆听的耳朵，请听我！我是您的仆人；今日我最谦卑地敬拜您，在您面前流下我的眼泪。您全然是我的天主；您是我的主，正如天使无数次告诉我的，特别是在那日，我的灵魂因对纯洁而蒙福的玛利亚——她正怀有您在她腹中，我曾想暗暗休弃她——起了悖逆的念头。我正如此思量时，看，在安息中主的天使向我显现，以奇妙的奥秘对我说：\InnerQuote{哦，若瑟，达味之子，不要害怕娶玛利亚为妻；不要使你的灵魂忧伤，也不要对她的怀孕说不合适的话，因为她是由圣神怀孕，将要生一个儿子，你要给他起名叫耶稣，因为祂要把自己的百姓从他们的罪恶中拯救出来。}求您不要因此以恶意待我，主啊！因为我当时对您诞生的奥秘一无所知。我的主，我也记得那日，一个孩子被蛇咬死。他的亲属想要把您交给黑落德，说您杀了他；但您使他从死者中复活，交还给了他们。于是我走到您跟前，抓住您的手说：\InnerQuote{我的儿子，您要当心。}但您回答我说：\InnerQuote{你不是我在肉身上的父亲吗？我要教导你我是谁。}所以如今，主和我的天主，请不要因那个时辰向我发怒或定我的罪。我是您的仆人，是您婢女的儿子；但您是我的主，我的天主和救主，最确实地是天主子。}\par

18. 当我的父亲若瑟这样说时，他不能再哭泣了。我看到死亡如今已掌管了他。我的母亲，无玷的童贞女，起身来到我面前说：\Quote{哦，我亲爱的儿子，这位虔诚的老人若瑟快要死了。}我回答说：\Quote{哦，我最亲爱的母亲，确实，世上一切受造物都同样面临死亡的必然性；因为死亡掌管全人类。连您，哦，我的童贞母亲，也必须像其他凡人一样等待生命的终结。然而您的死亡，以及这位虔诚之人的死亡，不是死亡，而是永存的生命。不仅如此，连我也必须死，就我从您所领受的身体而言。但请起来，哦，我可敬的母亲，进到若瑟——那位蒙福的老人那里去，好使您能看见他的灵魂离开身体时发生的事。}\par

19. 因此，我的无玷母亲玛利亚走去，进了若瑟所在的地方。我正坐在他脚边看着他，因为死亡的征兆已显现在他的面容上。那位蒙福的老人抬起头，眼睛一直盯着我的脸；但由于死亡的剧痛抓住了他，他无力对我说话。但他不断发出许多叹息。我握着他的手整整一个时辰；他把脸转向我，示意我不要离开他。此后，我把手放在他的胸口，感觉到他的灵魂已靠近喉咙，正预备离开它的寓所。\par

20. 当我的童贞母亲看见我抚摸他的身体时，她也触摸了他的脚。发现它们已经冰冷，没有了热气，她对我说：\Quote{我亲爱的儿子，他的脚确实已经开始僵硬，像雪一样冰冷。}于是她叫来他的儿女，对他们说：\Quote{你们所有的人都来，到你们的父亲那里去；因为他确实快要死了。}他的女儿阿西亚回答说：\Quote{哀哉，我的弟兄们，这确实是我亲爱的母亲得的那种病。}她痛哭流涕；若瑟的其他儿女都同她一起哀悼。我和我的母亲玛利亚也同他们一起哭泣。\par

21. 我转向南方，看见死亡已经临近，整个地狱与他一同前来，他的军队和随从紧紧跟随着他；他们的衣服、面容和口里都喷出火焰。当我的父亲若瑟看见他们径直向他走来时，他的眼睛泪水盈眶，同时发出奇怪的呻吟。当我看到他叹息的剧烈程度时，我击退了死亡和所有伴随他的仆从。我呼求我的善父，说：——\par

22. 一切仁慈的圣父，洞察的眼睛和垂听的耳朵，请垂听我为老人若瑟的祈祷和恳求；请派遣你的天使之长弥额尔，光明使者嘉俾厄尔，以及你所有光明的天使，让他们的全体行列陪伴我父亲若瑟的灵魂，直到他们将他引领到你面前。这是我父亲需要怜悯的时刻。我告诉你们：所有的圣人，以及世界上所有出生的人，无论是义人还是恶人，都必须尝到死亡。\par

23. 于是弥额尔和嘉俾厄尔来到我父亲若瑟的灵魂前，拿起它，用闪亮的包裹裹好。他就这样将灵魂交在我善父的手中，父赐给他平安。那时他的儿女还不知道他已经安息了。天使们保护他的灵魂不受路上黑暗的恶魔侵扰，赞美天主，直到他们将他领进虔诚者的居所。\par

24. 他的身体倒卧在地，毫无血色；因此我伸出手，合上他的眼睛，闭上他的嘴，对童贞玛利亚说：\Quote{我的母亲，他在世上生活时展现的技能哪里去了？看，它已经消失，好像从未存在过。}当他的儿女听见我和我的母亲，纯洁的童贞女说话时，他们知道他已断气，便流泪哀悼。但我对他们说：\Quote{你们父亲的死不是死亡，而是永生；因为他已脱离今生的烦恼，进入了永恒安息。}他们听了这些话，就撕裂衣服，哭泣起来。\par

25. 纳匝肋和加里肋亚的居民听见他们的哀哭，都聚集前来，从第三时辰哭到第九时辰。在第九时辰，他们一同来到若瑟的床前。他们抬起他的身体，用昂贵的香膏涂抹。但我用天上的祷文祈求我的圣父——就是在我还在童贞玛利亚、我的母亲腹中时亲手作的同一篇祷文。我刚念完并说阿们，就有大量天使前来；我命令其中两个天使展开他们闪亮的衣服，包裹若瑟——这位有福的老人——的身体。\par

26. 我对若瑟说：\Quote{死亡的腐朽气味不能统治你，虫也不能从你的身体里出来。你的肢体一根也不可折断，你头上的头发一根也不可改变。你的身体任何部分都不会朽坏，我的父亲若瑟，它将完整无缺、毫无腐朽地存留，直到千年盛宴。凡在你的纪念日献上祭品的人，我必祝福他，并在童女会中报答他；凡在自己的工作中、在你纪念的日子、以你的名将食物献给困苦、贫穷、寡妇和孤儿的人，他一生不会缺少美物。凡以你的名给寡妇或孤儿一杯水或酒喝的人，我必将他赐给你，使你与他一同进入千年盛宴。凡在你的纪念日献祭的人，我必祝福他，并在童女教会中报答他：我要以三十倍、六十倍、一百倍偿还他。凡记录你生平、你劳苦和你离世，以及从我口中发出的这篇叙述的人，我将他交给你保管，只要他尚在今生。当他的灵魂离开身体，必须离开这世界时，我要烧毁他的罪簿，在审判之日不以任何刑罚折磨他；他将越过火焰之海，毫无艰难痛苦。至于那些无法给出我所提及之物的穷人，他有义务：若他有一个儿子出生，就当给他起名叫若瑟。这样，那家里将永不发生贫穷或猝死。}\par

27. 此后，城中的首领们一同来到那蒙福老人\Person{Joseph}的遗体安放之处，带来了殓衣；他们想按犹太人惯于安葬亡者的方式，用殓衣包裹遗体。但他们发现他紧握着自己的裹尸布；因为那布紧紧粘着身体，当他们想取下时，竟如同铁一般，无法移动或松开。他们也找不到那块麻布的末端，这令他们极为惊异。最后，他们将遗体抬到一处有洞穴的地方，打开墓门，好将他的身体葬在他祖先的旁边。那时，我想起了他与我同往\Place{Egypt}的日子，以及他为我所受的极大痛苦。于是，我哀悼他的死良久，伏在他身上说道：——\par

28. \Quote{啊，死亡！你使一切知识消逝，引起如此多的眼泪和哀号，实在是我的父天主亲自赐予了你这种能力。因为人因\Person{Adam}和他妻子\Person{Eve}的过犯而死亡，死亡连一个人也不放过。然而，没有一件事临到任何人，或是加给他，是未经我父命令的。确实有人曾活到九百岁，但他们还是死了。是的，即使有些人活得更久，他们终究也遭遇了同样的命运；没有一个人曾说过：我没有尝过死亡。因为主从不施加同样的惩罚两次，既然我父乐意将它降在人身上。当它出发时，看见从天上降下的命令，它就说：我要出去攻击那人，大大地摇动他。然后它毫不迟延地攻击灵魂，取得对灵魂的掌控，任意对待它。因为\Person{Adam}没有遵行我父的旨意，反而违背了祂的诫命，我父的愤怒被激起，判他死亡；于是死亡进入了世界。但如果\Person{Adam}遵守了我父的规诫，死亡就绝不会临到他。你们以为我能请求我善父派一火战车来，\BibleRef{列下 2:11}载起我父\Person{Joseph}的身体，将他送到安息之所，好使他与神灵同住吗？但因\Person{Adam}的过犯，死亡的烦恼和暴力已降临到全人类。正因如此，我必须按肉体死去，为的是我创造的工程，好使他们获得恩宠。}\par

29. 我说完这些话后，拥抱了我父\Person{Joseph}的身体，为他哭泣；他们打开墓门，将他的身体安放在里面，靠近他父亲\Person{Jacob}的身体。在他安息时，他已活了一百一十一岁。他的口中从未有一颗牙齿疼痛，他的视力也未曾减退，身体没有佝偻，力气也未减弱；他一直以木匠手艺工作，直到生命的最后一天；那天是亚笔月二十六日。\par

30. \Quote{我们宗徒从救主那里听到这些事后，满心欢喜地起身，俯伏在地尊敬祂，说：我们的救主啊，请向我们显示祢的恩宠。如今我们确实听到了生命之言；然而，我们的救主啊，我们对\Person{Enoch}和\Person{Elias}的命运感到惊奇，因为他们不必经历死亡。确实，他们至今仍居住在义人的居所，他们的身体也未见腐朽。然而，那位年迈的木匠\Person{Joseph}，却是祢按肉身的父亲。祢曾命令我们往普天下去宣讲圣福音；祢曾说过：\InnerQuote{向他们讲述我父\Person{Joseph}的死，每年以隆重的礼仪为他举行庆节和圣日。凡从这记述中删去什么或添加什么的，就是犯罪。\BibleRef{默 22:18}}我们特别惊奇的是，\Person{Joseph}从祢在\Place{Bethlehem}出生的那一天起，就称祢为按肉身的儿子。那么，为何祢没有像对待他们那样使他成为不朽的呢？而祢又说他是义人和被拣选的？}\par

31. \Quote{确实，我父关于\Person{Adam}因悖逆而发的预言，如今已应验了。一切都按照我父的旨意和喜悦安排。因为若有人拒绝天主的诫命，跟随魔鬼的行为去犯罪，他的生命反而得以延长；他被保存是为了使他可能悔改，并想到他必将被交付于死亡之手。但若有人热心行善，他的生命也会延长，好使随着他年高德劭，正直的人可以效法他。但当你们看见一个人心地易怒，他的日子必会缩短；因为这些人正是在壮年时被取去的。因此，我父关于人类所说的一切预言，都必须逐字逐句地应验。至于\Person{Enoch}和\Person{Elias}，他们如何至今活着，保持着他们出生时的同一身体；以及关于我父\Person{Joseph}，他未能像他们一样留在身体里：确实，即使一个人活在世上千万年，终究还是得在某时换生命为死亡。我告诉你们，我的弟兄们，他们，\Person{Enoch}和\Person{Elias}，\BibleRef{默 11:3}在末世时必须返回世界并死去——就是在动荡、恐怖、困惑和痛苦的日子。因为假基督将杀死四个身体，将他们的血像水一样倒出，因为他们会揭露他的不敬，使他受到谴责，并在他们有生之年使他蒙受羞辱。}\par

32. 我们便说：\Quote{我主、天主和救主，祢所说的那四个人，祢说假基督将因他们加于他的羞辱而将他们剪除，他们是谁？}主回答说：\Quote{他们是\Person{厄诺客}、\Person{厄里亚}、\Person{希拉}和\Person{塔彼达}。}我们一听救主这话，便欣喜雀跃，将一切光荣与感谢归于主天主和我们的救主耶稣基督。光荣、尊威、尊严、统治、权能、赞美，都归于祂，并与祂一起归于良善的圣父，以及赋予生命的圣神，从今世直到永远。阿们。\par

\SourceRecord{Translated by Alexander Walker.；Ante-Nicene Fathers；Vol. 8.；Edited by Alexander Roberts, James Donaldson, and A. Cleveland Coxe.；Buffalo, NY: Christian Literature Publishing Co.,；1886.；<http://www.newadvent.org/fathers/0805.htm>.}
